\section{Implementation Details}
\label{sec:implementation}

\subsection{Code Organization}

Miniforge follows a modular package structure:

\begin{lstlisting}[language=bash, caption=Package Structure]
miniforge/
|-- __init__.py          # Main exports
|-- core/                # Core engine
|   |-- engine.py        # InferenceEngine
|   |-- memory.py        # MemoryManager
|   -- backends/        # Backend implementations
|       |-- llama_cpp.py
|       -- transformers.py
|-- models/              # Model management
|   |-- minimax.py       # Main interface
|   |-- registry.py      # Model registry
|   -- gguf_convert.py  # Conversion utilities
|-- generation/           # Generation utilities
|   |-- streaming.py
|   -- tools.py
|-- multimodal/          # Vision support
|   -- vision.py
-- utils/               # Utilities
    -- config.py        # M7Config
\end{lstlisting}

\subsection{Async Implementation}

Async programming patterns throughout Miniforge:

\subsubsection{Thread Pool Delegation}

CPU-intensive operations run in thread pools to avoid blocking:

\begin{lstlisting}[language=Python, caption=Thread Pool Pattern]
async def generate(self, prompt: str, **kwargs) -> str:
    if not self._initialized:
        await self.initialize()
    
    loop = asyncio.get_event_loop()
    result = await loop.run_in_executor(
        None,  # Default executor
        lambda: self._llm.create_completion(
            prompt=prompt, **kwargs
        )
    )
    return result["choices"][0]["text"]
\end{lstlisting}

\subsubsection{Streaming with Async Generators}

Token streaming uses async generators:

\begin{lstlisting}[language=Python, caption=Streaming Implementation]
async def generate_stream(self, prompt: str, **kwargs):
    loop = asyncio.get_event_loop()
    
    stream = await loop.run_in_executor(
        None,
        lambda: self._llm.create_completion(
            prompt=prompt, stream=True, **kwargs
        )
    )
    
    for chunk in stream:
        text = chunk["choices"][0].get("text", "")
        if text:
            yield text
        await asyncio.sleep(0)  # Yield control
\end{lstlisting}

\subsection{Configuration Management}

Configuration uses dataclasses with validation:

\begin{lstlisting}[language=Python, caption=Configuration Dataclass]
@dataclass
class M7Config:
    max_memory_gb: float = 24.0
    n_threads: int = 8
    n_ctx: int = 194_560
    quantization: str = "Q4_K_M"
    cache_type_k: str = "turbo3"
    cache_type_v: str = "turbo3"
    flash_attn: bool = True
    
    def __post_init__(self):
        # Validation
        valid_quants = ["Q2_K", "Q3_K_M", "Q4_K_M", "Q5_K_M", "Q6_K", "Q8_0"]
        if self.quantization not in valid_quants:
            self.quantization = "Q4_K_M"
\end{lstlisting}

\subsection{Error Handling}

Comprehensive error handling with graceful degradation:

\begin{lstlisting}[language=Python, caption=Error Handling]
try:
    from llama_cpp import Llama
except ImportError:
    raise ImportError(
        "llama-cpp-python not installed. "
        "Install with: uv pip install llama-cpp-python"
    )

# Context fallback
for candidate_ctx in _context_fallbacks(requested_n_ctx, max_safe_ctx):
    kwargs["n_ctx"] = candidate_ctx
    try:
        self._llm = await loop.run_in_executor(
            None, lambda: Llama(**kwargs)
        )
        break
    except ValueError as exc:
        logger.warning(f"Failed with n_ctx={candidate_ctx}: {exc}")
        continue
\end{lstlisting}

\subsection{Memory Management Implementation}

Memory tracking uses psutil for system monitoring:

\begin{lstlisting}[language=Python, caption=Memory Monitoring]
def get_stats(self) -> MemoryStats:
    mem = psutil.virtual_memory()
    return MemoryStats(
        total_gb=mem.total / (1024**3),
        available_gb=mem.available / (1024**3),
        used_gb=mem.used / (1024**3),
        percent_used=mem.percent,
        model_memory_gb=self.current_model_memory,
        kv_cache_gb=self.current_kv_memory,
    )
\end{lstlisting}

\subsection{Tool Calling Implementation}

Tool calling support with async execution:

\begin{lstlisting}[language=Python, caption=Tool Execution]
async def execute_single(self, tool_call: ToolCall) -> ToolResult:
    name = tool_call.name
    handler = self._handlers[name]
    
    try:
        if asyncio.iscoroutinefunction(handler):
            result = await handler(**tool_call.arguments)
        else:
            loop = asyncio.get_event_loop()
            result = await loop.run_in_executor(
                None,
                lambda: handler(**tool_call.arguments)
            )
        
        return ToolResult(
            tool_call_id=tool_call.id,
            content=str(result),
            is_error=False,
        )
    except Exception as e:
        return ToolResult(
            tool_call_id=tool_call.id,
            content=f"Error: {e}",
            is_error=True,
        )
\end{lstlisting}

\subsection{Vision Processing}

Vision support with PIL for image handling:

\begin{lstlisting}[language=Python, caption=Vision Processing]
class VisionProcessor:
    def preprocess(self, image, target_size=None):
        from PIL import Image
        
        if not isinstance(image, Image.Image):
            img = self.load_image(image)
        else:
            img = image
        
        # Convert to RGB
        if img.mode != "RGB":
            img = img.convert("RGB")
        
        # Resize
        if target_size:
            img = img.resize(target_size, Image.Resampling.LANCZOS)
        elif max(img.size) > self.max_size:
            img.thumbnail((self.max_size, self.max_size))
        
        return img
    
    def encode_base64(self, image, format="PNG") -> str:
        img = self.preprocess(image)
        buffer = BytesIO()
        img.save(buffer, format=format)
        return base64.b64encode(buffer.getvalue()).decode()
\end{lstlisting}

\subsection{Type Safety}

Strict typing throughout the codebase:

\begin{lstlisting}[language=Python, caption=Type Annotations]
from typing import Optional, Dict, Any, List, AsyncIterator, Union

async def chat(
    self,
    message: str,
    history: Optional[List[Dict[str, str]]] = None,
    system_prompt: Optional[str] = None,
    max_tokens: Optional[int] = None,
    temperature: Optional[float] = None,
    tools: Optional[List[Tool]] = None,
    stream: bool = False,
) -> Union[str, AsyncIterator[str]]:
    ...
\end{lstlisting}

\subsection{Testing Strategy}

Comprehensive test coverage:

\begin{itemize}
    \item \textbf{Unit tests:} Individual component testing
    \item \textbf{Integration tests:} End-to-end workflows
    \item \textbf{Benchmark tests:} Performance regression detection
    \item \textbf{Memory tests:} OOM prevention validation
\end{itemize}

\begin{lstlisting}[language=Python, caption=Example Test]
def test_memory_manager_select_quantization():
    mem = MemoryManager()
    
    # 2.7B model should fit with Q4_K_M
    quant = mem.select_quantization(2.7)
    assert quant in ["Q3_K_M", "Q4_K_M", "Q5_K_M", "Q6_K"]
    
    # 7B model might need more aggressive quantization
    quant_7b = mem.select_quantization(7.0)
    assert quant_7b in ["Q3_K_M", "Q4_K_M"]
\end{lstlisting}

\subsection{Documentation}

Documentation includes:

\begin{itemize}
    \item Docstrings for all public APIs
    \item Type hints for function signatures
    \item Usage examples in docstrings
    \item README with quickstart guide
    \item Architecture documentation
\end{itemize}
